\section{Conclusion, Limitation and Future Work}
In this work, we explore a scalable way for building a general representation model across different modalities.
Based on the flexible architecture and modality-agnostic pretraining tasks, we release \onepeace, a general representation model that can seamlessly align and integrate representations across vision, audio, and language modalities.
We conduct a series of experiments across 3 modalities, 11 tasks, and 16 datasets.
The experimental results demonstrate that \onepeace achieves leading results in a wide range of tasks, including image classification, semantic segmentation, audio-text retrieval, audio classification, audio question answering, image-text retrieval, and visual grounding.
Furthermore, we show that \onepeace possesses a strong emergent zero-shot retrieval capability, enabling it to align modalities that are not paired in the training data.

% In this work, we propose \onepeace, a general representation model that can seamlessly align and integrate representations across vision, audio, and language modalities.
% The architecture of \onepeace includes multiple modality adapters and a modality fusion encoder, which allow interaction between different modalities while facilitating information extraction within each modality.
% The contrastive-style pretraining tasks of \onepeace enable the model to achieve superior fine-tuning performance and allow for zero-shot cross-modal retrieval.
% We design two types of contrastive pretraining tasks: inter-modal contrastive learning and intra-modal contrastive learning with masking, with the aim of achieving outstanding finetuning performance while preserving the cross-modal zero-shot retrieval ability.
% To allow the model achieve great finetuning performance while maintaining cross-modal zero-shot retrieval ability, we design two types of contrastive-style pretraining tasks: inter-modal contrastive learning and masked intra-modal contrastive learning.
% We design two types of contrastive-style pretraining tasks: inter-modal contrastive learning and masked intra-modal contrastive learning. 
% These tasks can effectively align the semantic spaces of vision, audio and language.
% We conduct a series of experiments across 3 modalities, 11 tasks, and 16 datasets.
% The experimental results demonstrate that \onepeace achieves new SOTA results in a wide range of tasks, including image classification, semantic segmentation, audio-text retrieval, audio classification, audio question answering, image-text retrieval, and visual grounding.

\paragraph{Limitation.}
Although \onepeace achieves leading results in a wide range of tasks, it falls short of achieving state-of-the-art results in zero-shot image-text retrieval and vision-language understanding tasks.
There are two possible reasons for this:
1). \onepeace didn't see enough image-text pairs during pretraining. We only trained on 6.4 billion image-text pairs, while previous works~\cite{clip,align,coca} typically train on 12.8 billion image-text pairs or more.
2). \onepeace didn't use language pretrained models for initialization or introduce any pure text data. Both the vision and language modules of \onepeace are completely randomly initialized, while previous works~\cite{ofa,blip2} show that introducing pure text data or initialized with the language pretrained models can greatly enhance the model's performance. 
In fact, as a highly extensible model, \onepeace can combine with language pretrained models to achieve better results.

\paragraph{Future Work.}
In the future, we will test \onepeace on more downstream tasks, such as vision-audio-language tasks and extend to more modalities for pretraining like video, 3D point cloud, etc.
%Then, just two things:
Also, we pursue an active interaction with large language models (LLMs) to continue influencing broader areas. This includes:

\begin{itemize}
    % \item Build a more powerful general representation model with the help of large-scale language models;
    % \item Build a more general multimodal language model by combining large-scale language models.
    \item With the help of LLMs, building a more powerful general representation model.
    \item By combining LLMs, creating a more general multimodal language model.
\end{itemize}